% course.tex

%%%%%%%%%%%%%%%%%%%%
\begin{frame}{}
  \begin{center}
    \red{\bf \large 什么是``计算机数学''?}

    \pause
    \vspace{1em}
    \blue{\bf \large ``Mathematics for Computer Science''} \\[5pt]
    \teal{(\textsf{\bf math4cs})}

    \fig{width = 0.35\textwidth}{figs/math4cs-qq-logo}
  \end{center}
\end{frame}
%%%%%%%%%%%%%%%%%%%%

%%%%%%%%%%%%%%%%%%%%
\begin{frame}{}
  \begin{center}
    QQ 群号: \blue{\bf 108 745 6358} \qquad 验证问题答案: \teal{``\bf math4cs''}
    \fig{width = 0.40\textwidth}{figs/math4cs-qq-qrcode}
    修改群昵称: ``\teal{\bf 学号-姓名}''
  \end{center}
\end{frame}
%%%%%%%%%%%%%%%%%%%%

%%%%%%%%%%%%%%%%%%%%
\begin{frame}{}
  \begin{center}
    {\large ``一入\red{\bf 数学}深似海, 从此娱乐是路人''}
    \fig{width = 0.90\textwidth}{figs/math-map}
  \end{center}
\end{frame}
%%%%%%%%%%%%%%%%%%%%

%%%%%%%%%%%%%%%%%%%%
\begin{frame}{}
  \begin{center}
    \red{\bf \large What does ``Computer Science'' study?}

    \pause
    \vspace{2em}
    \begin{quote}
      Computer science focuses on methods involved in
      \blue{\underline{\bfit{specification}}, \underline{\bfit{design}},
      programming, \bfit{verification},
      implementation and testing} of \teal{human-made computing systems}.
    \end{quote}

    \pause
    \vspace{2em}
    \blue{\large math4cs is {\purple{\underline{\bfit{model-and-proof}}}} oriented.}
  \end{center}
\end{frame}
%%%%%%%%%%%%%%%%%%%%

%%%%%%%%%%%%%%%%%%%%
\begin{frame}{}
  \begin{center}
    {\bf \large \red{四大主题:} \teal{逻辑、集合论、图论、群论 (抽象代数)}}

    \fig{width = 0.50\textwidth}{figs/math-tree}

    \pause
    {\bf \blue{支流遍布:} \teal{组合与计数、\gray{数论}、\gray{(离散)概率}、$\ldots$}}

    \pause
    \vspace{1em}
    \red{\bf \large 计算机数学是个大杂烩, 啥都学点儿, \xout{啥都没学好}}
  \end{center}
\end{frame}
%%%%%%%%%%%%%%%%%%%%

% %%%%%%%%%%%%%%%%%%%%
% \begin{frame}{}
%   \begin{center}
%     \red{\bf \large ``计算机数学''在哪里?}
%     \fig{width = 0.60\textwidth}{figs/math-tree}

%     \pause
%     \blue{\bf \large 计算机数学是个大杂烩, 啥都学点儿, \xout{啥都没学好}}
%   \end{center}
% \end{frame}
% %%%%%%%%%%%%%%%%%%%%

%%%%%%%%%%%%%%%%%%%%
\begin{frame}{}
  \begin{center}
    \red{\bf \Large 计算机(离散)数学}

    \vspace{0.50cm}
    {\large 研究\blue{\bf 离散对象的结构、性质、操作等}的数学分支}
  \end{center}
\end{frame}
%%%%%%%%%%%%%%%%%%%%

%%%%%%%%%%%%%%%%%%%%
\begin{frame}{}
  \begin{center}
    \red{\Large 将计算机(离散)数学看作一门语言，一套工具}

    \vspace{1.00cm}
    {\large 培养形式化描述问题的能力}

    \vspace{0.60cm}
    {\large 培养做严格证明的能力}
  \end{center}
\end{frame}
%%%%%%%%%%%%%%%%%%%%

%%%%%%%%%%%%%%%%%%%%
\begin{frame}{}
  \begin{center}
    {\large 分班教学 (共 \blue{\bf 9} 个班级)}

    \vspace{0.30cm}
    {\large \red{拔尖班: 独立授课, 独立考核}}

    \pause
    \vspace{0.30cm}
    \fig{width = 0.35\textwidth}{figs/confident-student}
  \end{center}
\end{frame}
%%%%%%%%%%%%%%%%%%%%

%%%%%%%%%%%%%%%%%%%%
\begin{frame}{}
  \begin{center}
    {\large 平时作业 {\it vs.} 课堂测验 {\it vs.} 期中测试 {\it vs.} 期末测试}

    \[
      1 \quad:\quad 3 \quad:\quad 2 \quad:\quad 4
    \]
  \end{center}
\end{frame}
%%%%%%%%%%%%%%%%%%%%

%%%%%%%%%%%%%%%%%%%%
\begin{frame}{}
  \begin{center}
    \red{\bf 每周四}下午 14:00 发布平时作业 \qquad \red{\bf 下周四}晚 22:00 前提交作业

    \vspace{1em}
    取最高的 \blue{\bf 12} 次作业成绩计入总评

    \vspace{1em}
    \fig{width = 0.30\textwidth}{figs/do-math}

    \red{\large \bf 请按时提交, 过时不补}
  \end{center}
\end{frame}
%%%%%%%%%%%%%%%%%%%%

%%%%%%%%%%%%%%%%%%%%
\begin{frame}{}
  \begin{center}
    \teal{``头歌 (EduCoder)''课程邀请码: AF7LD}

    \begin{columns}
      \column{0.50\textwidth}
        \fig{width = 0.80\textwidth}{figs/educoder-qrcode}
      \column{0.50\textwidth}
        \pause
        \fig{width = 0.30\textwidth}{figs/tex}
        \vspace{-0.30cm}
        \fig{width = 1.00\textwidth}{figs/math4cs-problem-sets-github}
    \end{columns}

    \vspace{0.50cm}
    \teal{\small \url{https://github.com/courses-at-hnu-by-hfwei/math4cs-problem-sets}}
  \end{center}
\end{frame}
%%%%%%%%%%%%%%%%%%%%

%%%%%%%%%%%%%%%%%%%%
\begin{frame}{}
  \begin{center}
    约\red{\bf 每两周一次}课堂测验 (提前通知时间与测验范围)

    \vspace{1em}
    取最高的 \blue{\bf $5 \sim 6$} 次课堂测验成绩计入总评

    \vspace{1em}
    \fig{width = 0.30\textwidth}{figs/academic-fortress}

    \red{\large \bf 请准时参加, 不安排补考}
  \end{center}
\end{frame}
%%%%%%%%%%%%%%%%%%%%

%%%%%%%%%%%%%%%%%%%%
\begin{frame}{}
  \fig{width = 0.50\textwidth}{figs/yuefa}
\end{frame}
%%%%%%%%%%%%%%%%%%%%

%%%%%%%%%%%%%%%%%%%%
\begin{frame}{}
  \begin{center}
    \red{\Large \bf 非必要, 不点名}
  \end{center}
\end{frame}
%%%%%%%%%%%%%%%%%%%%

%%%%%%%%%%%%%%%%%%%%
\begin{frame}{}
  \begin{center}
    \red{\bf \Large 非必要, 不迟到}

    \pause
    \vspace{0.60cm}
    \blue{\bf \large 尽量吃早餐, 但不可以在教室吃早餐}
  \end{center}
\end{frame}
%%%%%%%%%%%%%%%%%%%%

%%%%%%%%%%%%%%%%%%%%
\begin{frame}{}
  \begin{center}
    \red{\bf \Large \xout{非必要}, 不抄袭; 一经发现, 后果严重}

    \pause
    \vspace{1.50cm}
    \blue{\bf \Large {当次作业计 0 分;} 平时作业成绩扣 2 分, 扣完为止}
  \end{center}
\end{frame}
%%%%%%%%%%%%%%%%%%%%

%%%%%%%%%%%%%%%%%%%%
\begin{frame}{}
  \begin{center}
    \fig{width = 0.45\textwidth}{figs/discrete-math-book-yang}

    \vspace{1em}
    {\large 授课内容不局限于教材, 认真听讲很重要}
  \end{center}
\end{frame}
%%%%%%%%%%%%%%%%%%%%

%%%%%%%%%%%%%%%%%%%%
\begin{frame}{}
  \begin{columns}
    \column{0.50\textwidth}
      \fig{width = 0.80\textwidth}{figs/dm-structures-ch}
    \column{0.50\textwidth}
      \fig{width = 0.80\textwidth}{figs/dm-applications-book}
  \end{columns}

  \begin{center}
    {\large 内容与习题偏简单, 略显琐碎; 可用于课前预习及课后基础练习}
  \end{center}
\end{frame}
%%%%%%%%%%%%%%%%%%%%

%%%%%%%%%%%%%%%%%%%%
\begin{frame}{}
  \begin{columns}
    \column{0.50\textwidth}
      \fig{width = 0.80\textwidth}{figs/dm-book-red}
    \column{0.50\textwidth}
      \fig{width = 0.80\textwidth}{figs/dm-book-green}
  \end{columns}

  \begin{center}
    {\large 带详细答案的习题集; 指定部分习题供课后观摩}
  \end{center}
\end{frame}
%%%%%%%%%%%%%%%%%%%%

%%%%%%%%%%%%%%%%%%%%
\begin{frame}{}
  \begin{center}
    \fig{width = 0.50\textwidth}{figs/mcs-book}

    \vspace{1em}
    {\large \blue{\bf 推荐阅读}; 其它参考书随课程进度、按专题指定阅读}
  \end{center}
\end{frame}
%%%%%%%%%%%%%%%%%%%%

%%%%%%%%%%%%%%%%%%%%
\begin{frame}{}
  \begin{center}
    \fig{width = 0.80\textwidth}{figs/math4cs-lectures-github}

    \vspace{1em}
    \teal{\url{https://github.com/courses-at-hnu-by-hfwei/math4cs-lectures}}

    \vspace{2em}
    {0-overview.pdf \qquad 0-overview-\red{\bf handout}.pdf}
  \end{center}
\end{frame}
%%%%%%%%%%%%%%%%%%%%